% Generated draft: requires mathematical and source review.
\documentclass[11pt]{article}
\usepackage[margin=25mm]{geometry}
\usepackage{fontspec}
\setmainfont{Latin Modern Roman}
\usepackage{amsmath,amssymb,mathtools,unicode-math}
\setmathfont{Latin Modern Math}
\usepackage{enumitem,needspace,xurl,hyperref}
\setlength{\parindent}{0pt}
\setlength{\parskip}{5pt}
\setlength{\emergencystretch}{3em}
\newcommand{\N}{\mathbb{N}}
\newcommand{\Z}{\mathbb{Z}}
\newcommand{\Q}{\mathbb{Q}}
\newcommand{\R}{\mathbb{R}}
\newcommand{\C}{\mathbb{C}}
\newcommand{\F}{\mathbb{F}}
\newcommand{\A}{\mathbb{A}}
\newcommand{\PP}{\mathbb{P}}
\newcommand{\E}{\mathbb{E}}
\newcommand{\eps}{\varepsilon}
\newcommand{\dd}{\,\mathrm{d}}
\begin{document}
\section*{UnsolvedMath ID: 20000450}
\textbf{Source problem code:} catalogue entry
\subsection{Definitions and mathematical statement}
Let \(P\) denote the product of the five side-line equations of a regular pentagon in an affine plane, and \(C\) the equation of its circumcircle. McCallum's question asks to compute the five-torsion associated with
\[
P+\lambda C^2=0.
\]
This is an affine degree-five equation; if \(z\) homogenizes the plane, its projective form is \(P_h+\lambda zC_h^2=0\). For parameters where the only singularities are the stated five ordinary double points and the normalization is connected, the genus is one. Those restrictions and the coefficient field are not specified in the source. Five points lie at infinity, determined by the side directions.

For a genus-one normalization \(C_\lambda\), five-torsion is intrinsically \(J_\lambda[5]\) for \(J_\lambda=\operatorname{Jac}(C_\lambda)\). Identifying the curve itself with an elliptic curve requires an origin, possibly after extension. The attached aid reconstructs a fixed normalized regular-pentagon model over \(k=\Q(\sqrt5)\). For its generic parameter, or admissible specializations under the stated model hypotheses, it reports that the infinity points become
\[
O,D,2D,3D,4D
\]
after choosing \(O\), with \(D\) of exact order five. This lists one cyclic line, not all twenty-five geometric five-torsion points.

Put \(a=5+2\sqrt5\), and let \(\chi_a\) be the quadratic character of \(k(\sqrt a)/k\), extended to the model's parameter field when necessary. With \(\chi_5\) the mod-five cyclotomic character, the aid reports a canonical exact sequence of Galois modules
\[
0\longrightarrow\F_5(\chi_a)\longrightarrow J_\lambda[5]
\longrightarrow\F_5(\chi_5\chi_a)\longrightarrow0.
\]
The quotient character follows from the Weil pairing, and the aid explicitly does not assert that the sequence splits. The origin's field of definition and torsor descent must be distinguished from the Jacobian representation.

It also reports that reflection \(Y\mapsto-Y\) is negation after the appropriate origin choice, and that exact elimination plus a square-discriminant identity normalizes its quotient to a rational conic. This gives a one-variable double-cover reduction for the remaining Weierstrass and full division-field computation, but the conic equation, identity, original normalized model and extension class are not supplied here.

The source remarks motivate these genus-one curves as principal homogeneous spaces and candidates for five-torsion in a group extracted as ``X,'' and suggest a twist of the universal family over \(X_1(5)\). They also ask how the pencil relates to the conference talks, whether regularity of the pentagon is needed, and about variants using a star pentagon. These broader torsor/modular and variation questions are not settled by computing the single infinity line.
\subsection{Short English statement}
Compute five-torsion for the genus-one normalization of the pentagon/circle quintic. The attached aid reports the infinity cyclic subgroup and its Galois character, plus a nonsplit-allowed exact sequence and quotient reduction; the full division field and broader torsor questions remain unfinished.
\subsection{Sources}
\begin{itemize}
\item Local UnsolvedMath registry, permanent ID \(20000450\), exact public code \texttt{catalogue entry}; canonical location \texttt{aim-algebraic-number-theory-notes.json}, notes[101].
\item AIM workshop ``Rational and integral points on higher dimensional varieties,'' Question 17, attributed to McCallum; supplied URL \url{https://aimath.org/WWN/qptsurface2/qptsurface2.pdf}. Remarks are attributed to McCallum, Ellenberg and Voloch, mentioning \(X_1(5)\), principal homogeneous spaces and star-pentagon variants. No year is supplied; the literature field is empty. The PDF was not consulted here.
\item Attached research attempt 1, ``The infinity 5-torsion line and Kummer quotient of the pentagonal quintic,'' without full bibliographic references or its explicit normalized/conic equations.
\end{itemize}
\end{document}
